\chapter{处理—混杂反馈}

上一章把序贯可交换性确立为识别时变处理之因果效应的关键条件。设想我们有一项研究，其中最强形式的序贯可交换性成立：所测的时变混杂因子足以有效估计任何处理策略的因果效应。那么接下来的问题就是：该用哪种混杂校正方法。这个问题的答案，凸显了时变处理因果推断中的一个关键问题：处理—混杂反馈。

存在处理—混杂反馈时，使用传统校正方法可能使效应估计产生偏倚。\mnote{也就是说，即便我们掌握了有效估计任何处理策略之平均因果效应所需的全部信息，一般也无法做到。}本章描述处理—混杂反馈的结构，以及传统校正方法失效的原因。

\tikzset{
  nd/.style={inner sep=2pt,font=\small},
  bx/.style={draw,thick,inner sep=3.5pt,font=\small},
}

\section{处理—混杂反馈的要素}

再考虑上一章讨论过的 HIV 感染者序贯随机化试验。对研究中的每个人，我们在随访的每个月 $k=0,1,2\ldots K$ 都有处理 $A_k$（1：受处理，0：未处理）与协变量 $L_k$ 的数据，以及一个度量第 $K+1$ 个月健康状况的结局 $Y$。图~20.1 的因果图与图~19.2 相同，刻画了研究的前两个月。时变协变量 $L_k$ 是时变混杂因子。（与上一章一样，这里使用无删失的例子，以便专注于混杂。）

\begin{figure}[htbp]
\centering
\begin{tikzpicture}[causal,every node/.style={transform shape},scale=0.95]
  \node[nd] (L0) at (0,0) {$L_0$};
  \node[nd] (A0) at (2.2,0) {$A_0$};
  \node[nd] (L1) at (4.4,0) {$L_1$};
  \node[nd] (A1) at (6.3,0) {$A_1$};
  \node[nd] (Y)  at (8.0,0) {$Y$};
  \node[nd] (U0) at (0,-2.0) {$U_0$};
  \node[nd] (U1) at (4.4,-2.0) {$U_1$};
  \draw[ed] (L0) -- (A0);
  \draw[ed] (L0) .. controls (1.1,1.35) and (3.3,1.35) .. (L1);
  \draw[ed] (L0) .. controls (1.8,2.45) and (5.2,2.45) .. (A1);
  \draw[ed] (L0) .. controls (2.4,3.00) and (6.2,3.00) .. (Y);
  \draw[ed] (A0) -- (L1);
  \draw[ed] (A0) .. controls (3.3,1.05) and (5.4,1.05) .. (A1);
  \draw[ed] (A0) .. controls (4.2,1.72) and (6.6,1.72) .. (Y);
  \draw[ed] (A0) -- (U1);
  \draw[ed] (L1) -- (A1);
  \draw[ed] (L1) .. controls (5.7,0.80) and (7.2,0.75) .. (Y);
  \draw[ed] (A1) -- (Y);
  \draw[ed] (U0) -- (L0);
  \draw[ed] (U0) .. controls (1.6,-1.60) and (3.2,-1.05) .. (L1);
  \draw[ed] (U0) -- (U1);
  \draw[ed] (U0) .. controls (2.8,-3.40) and (6.8,-2.80) .. (Y);
  \draw[ed] (U1) -- (L1);
  \draw[ed] (U1) .. controls (5.9,-2.50) and (7.2,-1.40) .. (Y);
\end{tikzpicture}
\caption{时变混杂与处理—混杂反馈（两时点研究）}
\label{fig:20.1}
\end{figure}

图~20.1 中还有别的名堂。不仅有从 CD4 细胞计数 $L_k$ 指向处理 $A_k$ 的箭头，还有从处理 $A_{k-1}$ 指向未来 CD4 细胞计数 $L_k$ 的箭头——因为接受处理 $A_{k-1}$ 会提高未来的 CD4 细胞计数 $L_k$。也就是说，混杂因子影响处理，处理又影响混杂因子。这就是处理—混杂反馈（另见精要~20.1）。

\begin{finepoint}[用无环图表示反馈环]
有意思的是，像图~20.1 这样的无环图也能表示处理—混杂反馈的环路。做到这一视觉表示的诀窍，是把处理—混杂反馈的环路在时间上展开，即 $A_{k-1}\to L_k\to A_k\to L_{k+1}$，依此类推。

用无环图表示反馈环，还要求把时间视为离散变量。也就是说，我们说处理与协变量可以在每个区间 $[k,k+1)$ 内变化，$k=0,1,\ldots K$，但不指明变化发生在区间内的哪个确切时刻。这种时间离散化在实践中并非限制：区间的长度可以取得任意短，只要数据的粒度允许。例如，在一项个体每月或更少频次看医生的研究（如我们的 HIV 例子）中，把时间安全地离散为月区间即可。在其他情形下，年区间或日区间可能更合适。此外，正如我们在第 17 章所说，时间通常本就以离散区间（年、月、日）测量，因此时间的离散化往往根本算不上一种选择。
\end{finepoint}

注意，时变混杂可以在没有处理—混杂反馈的情况下出现。图~20.2 的因果图与图~20.1 相同，只是删去了从处理 $A_{k-1}$ 指向未来 $L_k$ 与 $U_k$ 的箭头。在由图~20.2 刻画的情形中，时变协变量 $L_k$ 是时变混杂因子，但它们不受先前处理的影响。因此，存在时变混杂，却不存在处理—混杂反馈。

\begin{figure}[htbp]
\centering
\begin{tikzpicture}[causal,every node/.style={transform shape},scale=0.95]
  \node[nd] (L0) at (0,0) {$L_0$};
  \node[nd] (A0) at (2.2,0) {$A_0$};
  \node[nd] (L1) at (4.4,0) {$L_1$};
  \node[nd] (A1) at (6.3,0) {$A_1$};
  \node[nd] (Y)  at (8.0,0) {$Y$};
  \node[nd] (U0) at (0,-2.0) {$U_0$};
  \node[nd] (U1) at (4.4,-2.0) {$U_1$};
  \draw[ed] (L0) -- (A0);
  \draw[ed] (L0) .. controls (1.1,1.35) and (3.3,1.35) .. (L1);
  \draw[ed] (L0) .. controls (1.8,2.45) and (5.2,2.45) .. (A1);
  \draw[ed] (L0) .. controls (2.4,3.00) and (6.2,3.00) .. (Y);
  \draw[ed] (A0) .. controls (3.3,1.05) and (5.4,1.05) .. (A1);
  \draw[ed] (A0) .. controls (4.2,1.72) and (6.6,1.72) .. (Y);
  \draw[ed] (L1) -- (A1);
  \draw[ed] (L1) .. controls (5.7,0.80) and (7.2,0.75) .. (Y);
  \draw[ed] (A1) -- (Y);
  \draw[ed] (U0) -- (L0);
  \draw[ed] (U0) .. controls (1.6,-1.60) and (3.2,-1.05) .. (L1);
  \draw[ed] (U0) -- (U1);
  \draw[ed] (U0) .. controls (2.8,-3.40) and (6.8,-2.80) .. (Y);
  \draw[ed] (U1) -- (L1);
  \draw[ed] (U1) .. controls (5.9,-2.50) and (7.2,-1.40) .. (Y);
\end{tikzpicture}
\caption{有时变混杂、无处理—混杂反馈}
\label{fig:20.2}
\end{figure}

处理—混杂反馈给因果推断出了一个有意思的难题。为把问题讲得尽量简单，我们再把图~20.1 的因果图简化一点。图~20.3 是图~20.1 中能说明两时点序贯随机化试验里处理—混杂反馈的最小子集。画图~20.3 时，我们作了四点简化：

\begin{itemize}
\item 由于我们关心的是 $L_1$ 所致混杂的后果，就不必为基线 CD4 细胞计数加入节点 $L_0$。只需假定处理 $A_0$ 是边际随机化的，处理 $A_1$ 在给定 $L_1$ 的条件下随机化。
\item 不纳入未测量的变量 $U_0$。
\item 没有从 $A_0$ 指向 $A_1$ 的箭头，这意味着处理的分配只用到 $L_1$ 的信息。
\item 没有从 $A_0$、$L_1$ 与 $A_1$ 指向 $Y$ 的箭头；若处理对任何个体的结局 $Y$ 都没有因果效应，即严格因果零假设成立，就会出现这种情形。
\end{itemize}

这些简化都不影响下面的论证。更复杂的因果图不会给本章的讨论增添任何概念上的洞见，只会更难读。

\begin{figure}[htbp]
\centering
\begin{tikzpicture}[causal,every node/.style={transform shape},scale=1.15]
  \node[nd] (A0) at (0,0) {$A_0$};
  \node[nd] (L1) at (2.3,0) {$L_1$};
  \node[nd] (A1) at (4.1,0) {$A_1$};
  \node[nd] (Y)  at (5.8,0) {$Y$};
  \node[nd] (U1) at (2.3,-1.5) {$U_1$};
  \draw[ed] (A0) -- (L1);
  \draw[ed] (L1) -- (A1);
  \draw[ed] (U1) -- (L1);
  \draw[ed] (U1) .. controls (4.5,-1.15) and (5.5,-0.75) .. (Y);
\end{tikzpicture}
\caption{两时点序贯随机化试验中处理—混杂反馈的最简因果图}
\label{fig:20.3}
\end{figure}

现在假定处理对任何个体的 $Y$ 都没有效应——这意味着图~20.3 的因果图是正确的，但研究者并不知道这一点。再假定我们有第 0 个月的处理 $A_0$、第 1 个月的处理 $A_1$、第 1 个月开始时的混杂因子 CD4 细胞计数 $L_1$，以及随访结束时的结局 $Y$ 的数据。我们希望用这些数据估计静态处理策略“始终治疗”（$a_0=1,a_1=1$）相对于静态处理策略“从不治疗”（$a_0=0,a_1=0$）对结局 $Y$ 的平均因果效应，即
\[
\E[Y^{a_0=1,a_1=1}]-\E[Y^{a_0=0,a_1=0}]
\]
按照图~20.3，真实但研究者未知的平均因果效应为 0，因为从任一处理变量出发都没有前向路径通向结局。也就是说，无论从 $A_0$ 还是 $A_1$ 出发，顺着箭头的方向走，都到不了 $Y$。

图~20.3 可以刻画一项序贯随机化试验，因为没有从未测量的 $U$ 直接指向处理变量的箭头。因此，如上一章所述，我们应当能够用 $A_0$、$L_1$、$A_1$ 与 $Y$ 的观测数据得出 $\E[Y^{a_0=1,a_1=1}]-\E[Y^{a_0=0,a_1=0}]$ 等于 0 的结论。然而，正如下一节将要说明的，当我们用分层、结局回归、匹配等传统方法校正 $L_1$ 时，一般无法正确估计因果效应。也就是说，在这个例子里，试图用这些方法校正混杂因子 $L_1$，一般会得到一个不等于 0 的效应估计，因而无效。

\section{传统方法的偏倚}

换言之，存在时变混杂因子与处理—混杂反馈时，传统方法不能用来正确校正这些混杂因子。即便我们有足够的纵向数据来保证序贯可交换性，传统方法一般也不能给出任何处理策略之因果效应的有效估计。相反，g 方法即使在存在处理—混杂反馈时，也能适当地校正时变混杂因子。\mnote{图~20.3 既可以表示序贯随机化试验，也可以表示没有未测量混杂的观察性研究；图~20.4 表示观察性研究。}

传统方法的这一局限，适用于时变混杂因子受先前处理影响的情形（如图~20.3），也适用于时变混杂因子与先前处理共享原因 $W$ 的情形（如图~20.4，它是图~19.4 的子集）。我们把图~20.3 与图~20.4（以及图~19.2 与图~19.4）都称为处理—混杂反馈的例子。下一节解释传统方法为何无法恰当处理处理—混杂反馈。

\begin{figure}[htbp]
\centering
\begin{tikzpicture}[causal,every node/.style={transform shape},scale=1.15]
  \node[nd] (A0) at (0,0) {$A_0$};
  \node[nd] (L1) at (2.5,0) {$L_1$};
  \node[nd] (A1) at (4.3,0) {$A_1$};
  \node[nd] (Y)  at (6.0,0) {$Y$};
  \node[nd] (W0) at (0,-1.5) {$W_0$};
  \node[nd] (U1) at (2.5,-1.5) {$U_1$};
  \draw[ed] (W0) -- (A0);
  \draw[ed] (W0) -- (L1);
  \draw[ed] (L1) -- (A1);
  \draw[ed] (U1) -- (L1);
  \draw[ed] (U1) .. controls (4.7,-1.15) and (5.7,-0.75) .. (Y);
\end{tikzpicture}
\caption{时变混杂因子与先前处理共享未测量原因 $W_0$ 的观察性研究}
\label{fig:20.4}
\end{figure}

为说明传统方法的偏倚，考虑一项（假想的）序贯随机化试验：32\,000 名 HIV 感染者，两个时点 $k=0$ 与 $k=1$。处理 $A_0=1$ 在基线以 0.5 的概率随机分配。处理 $A_1$ 在第 1 个月随机分配，其概率只取决于第 1 个月开始时 CD4 细胞计数 $L_1$ 的取值——若 $L_1=0$（高），为 0.4；若 $L_1=1$（低），为 0.8。\mnote{这是一项理想试验：完全依从所分配的处理策略，且无失访。}结局 $Y$ 在随访结束时测量，它是 CD4 细胞计数、血清中病毒浓度及其他临床指标的函数，$Y$ 取值越高表示健康状况越好。

表~20.1 给出该试验的数据。为节省篇幅，表中每个 $A_0$、$L_1$、$A_1$ 取值组合占一行，而不是每个个体占一行。对八种组合中的每一种，表给出受试者人数 $N$ 与结局均值 $\E[Y\given A_0,L_1,A_1]$。于是第 1 行表明，2400 名 $(A_0=0,L_1=0,A_1=0)$ 个体的结局均值为 $\E[Y\given A_0=0,L_1=0,A_1=0]=84$。在这项序贯随机化试验中，可识别性条件——序贯可交换性、正值性、一致性——都成立。按设计，$A_0$ 对 $Y$ 的效应没有混杂因子，而 $L_1$ 是 $A_1$ 对 $Y$ 之效应的唯一混杂因子，因此（在给定 $L_1$ 的条件下）序贯可交换性成立。\mnote{若在更多的时点 $k$ 上处理 $A_k$ 受 $L_k$ 影响，那么 $L_k$ 就是时变混杂因子。}查看表~20.1 即可断定正值性条件满足，否则八行中就会有一行或几行的人数为零。

\begin{table}[htbp]
\centering
\caption{序贯随机化试验的数据：每种 $(A_0,L_1,A_1)$ 组合的人数与结局均值}
\label{tab:20.1}
\small
\begin{tabular}{ccccc}
\toprule
$N$ & $A_0$ & $L_1$ & $A_1$ & 结局均值 $\E[Y\given A_0,L_1,A_1]$ \\
\midrule
2400 & 0 & 0 & 0 & 84 \\
1600 & 0 & 0 & 1 & 84 \\
2400 & 0 & 1 & 0 & 52 \\
9600 & 0 & 1 & 1 & 52 \\
4800 & 1 & 0 & 0 & 76 \\
3200 & 1 & 0 & 1 & 76 \\
1600 & 1 & 1 & 0 & 44 \\
6400 & 1 & 1 & 1 & 44 \\
\bottomrule
\end{tabular}
\end{table}

图~20.3 的因果图刻画的是严格因果零假设成立时的这项序贯随机化实验。为检验表~20.1 的数据是否与图~20.3 的因果图相容，我们可以在过去协变量与处理的各水平内，分别估计各时点固定处理 $A_0$ 与 $A_1$ 的平均因果效应，它们都应当为零。下面的计算忽略随机变异。\mnote{图~20.3 刻画的是零假设，因为没有从 $L_1$ 指向 $Y$ 的箭头。否则 $A_0$ 会通过 $L_1$ 对 $Y$ 有效应。}

粗看表即可发现，处理 $A_1$ 的平均因果效应在 $A_0$ 与 $L_1$ 定义的全部四个层中确实为零。考虑 $A_0=0$、$L_1=0$ 的 4000 名个体中 $A_1$ 的效应，其数据见表~20.1 的第 1、2 行。第 1 期未受处理者的结局均值 $\E[Y\given A_0=0,L_1=0,A_1=0]$ 为 84，第 1 期受处理者的结局均值 $\E[Y\given A_0=0,L_1=0,A_1=1]$ 也是 84。因此差值
\[
\E[Y\given A_0=0,L_1=0,A_1=1]-\E[Y\given A_0=0,L_1=0,A_1=0]
\]
为零。由于可识别性条件成立，这一关联差能有效估计层 $(A_0=0,L_1=0)$ 中的平均因果效应
\[
\E\bigl[Y^{a_1=1}\given A_0=0,L_1=0\bigr]-\E\bigl[Y^{a_1=0}\given A_0=0,L_1=0\bigr]
\]
同样容易检验，处理 $A_1$ 对 $Y$ 的平均因果效应在其余三层 $(A_0=0,L_1=1)$、$(A_0=1,L_1=0)$、$(A_0=1,L_1=1)$ 中也为零，只需分别比较第 3 与第 4 行、第 5 与第 6 行、第 7 与第 8 行的结局均值。

现在可以证明 $A_0$ 的平均因果效应也为零。为此需要计算关联差 $\E[Y\given A_0=1]-\E[Y\given A_0=0]$；由于随机化，它是因果对比 $\E[Y^{a_0=1}]-\E[Y^{a_0=0}]$ 的有效估计量。第 0 期未受处理的 16\,000 名个体中，结局均值 $\E[Y\given A_0=0]$ 是第 1、2、3、4 行结局均值的加权平均，等于 60。\mnote{\parbox[t]{3.0cm}{\raggedright\scriptsize $\frac{2400}{16000}\cdot 84+\frac{1600}{16000}\cdot 84$\\[2pt] ${}+\frac{2400}{16000}\cdot 52+\frac{9600}{16000}\cdot 52$\\[2pt] ${}=60$}}类似算得的 $\E[Y\given A_0=1]$ 也是 60。因此 $A_0$ 的平均因果效应为零。

\begin{techpoint}[g 零检验]
设严格因果零假设为真，则任何反事实结局 $Y^{g}$ 都等于观测结局 $Y$。在这种情形下，对一切 $Y^{g}$ 的序贯可交换性，在两时点的研究中可以写作 $Y\indep A_0\given L_0$ 与 $Y\indep A_1\given A_0,L_0,L_1$。（这里用到了这样一个事实：对任意 $a_0$ 与 $l_0$ 的取值，都存在满足 $a_0=g(l_0)$ 的策略 $g$。）因此，在序贯可交换性下，检验这些条件独立关系就是检验严格因果零假设。这就是 g 零检验（Robins 1986）。注意，第一个独立关系意味着在 $L_0$ 定义的任何层中 $A_0$ 都没有因果效应，第二个独立关系意味着在 $L_1$ 与 $A_0$ 定义的任何层中 $A_1$ 都没有因果效应。

更一般地，Robins (1986) 的 g 零定理指出：在对一切 $g$ 的序贯随机化下，上述两个独立关系成立，当且仅当 $Y^{g}$ 的分布、因而均值 $\E[Y^{g}]$ 对一切 $g$ 都相同，且等于观测到的 $Y$ 的分布与均值。
\end{techpoint}

我们已经确认：把 $A_0$ 与 $A_1$ 分别当作时点固定处理时，它们（在给定过去的条件下）的因果效应为零。那么，若把联合处理 $(A_0,A_1)$ 当作时变处理、比较两种处理策略，结果又如何？例如，设我们想比较“始终治疗”与“从不治疗”，即 $(a_0=1,a_1=1)$ 对 $(a_0=0,a_1=0)$。由于可识别性条件成立，表~20.1 的数据应当足以有效估计这一效应。

由于策略中每个成分 $a_0$ 与 $a_1$ 的效应均为零，由 g 零定理（见技术注记~20.1）可知，平均因果效应 $\E[Y^{a_0=1,a_1=1}]-\E[Y^{a_0=0,a_1=0}]$ 为零。但是，若我们使用常规分析方法，从数据中得出的结论是这样吗？为回答这个问题，我们做两项数据分析。第一项不校正混杂因子 $L_1$，这应当给出错误的效应估计。第二项通过分层校正混杂因子 $L_1$。

\begin{enumerate}
\item 我们把两个时点都受处理的 9600 名个体（表~20.1 第 6、8 行）的结局均值，与两个时点都未受处理的 4800 名个体（第 1、3 行）相比。两个均值分别为 $\E[Y\given A_0=1,A_1=1]=54.7$ 与 $\E[Y\given A_0=0,A_1=0]=68$。\mnote{\parbox[t]{3.0cm}{\raggedright\scriptsize $\frac{3200}{9600}\cdot 76+\frac{6400}{9600}\cdot 44$\\[2pt] ${}=54.7$\\[5pt] $\frac{2400}{4800}\cdot 84+\frac{2400}{4800}\cdot 52$\\[2pt] ${}=68.0$}}关联差为 $54.7-68=-13.3$；若作因果解释，就意味着两个时点都不受处理比两个时点都受处理更好。这一分析给出了错误答案——一个非零的差值——因为 $\E[Y\given A_0=a_0,A_1=a_1]$ 不是 $\E[Y^{a_0,a_1}]$ 的有效估计量。需要校正混杂因子 $L_1$。

\item 我们通过分层校正 $L_1$。也就是说，在 $L_1$ 的各水平内，比较两个时点都受处理者与两个时点都未受处理者的结局均值。例如取层 $L_1=0$。两个时点都受处理者的结局均值 $\E[Y\given A_0=1,L_1=0,A_1=1]$ 为 76（第 6 行）。两个时点都未受处理者的结局均值 $\E[Y\given A_0=0,L_1=0,A_1=0]$ 为 84（第 1 行）。关联差为 $76-84=-8$；若作因果解释，就意味着在层 $L_1=0$ 中，两个时点都不受处理比两个时点都受处理更好。类似地，层 $L_1=1$ 中的差值 $\E[Y\given A_0=1,L_1=1,A_1=1]-\E[Y\given A_0=0,L_1=1,A_1=0]$ 也是 $-8$。
\end{enumerate}

我们说过效应估计应当是 0，而不是 $-8$。

怎么回事？校正了混杂因子的分析，怎么会也给出错误答案？这个估计反映了存在处理—混杂反馈时传统校正方法的偏倚。\mnote{注意：既然在 $L_1$ 的两个层中效应都是 $-8$，层特异效应的任何加权平均都不可能给出正确的 0。}下一节解释偏倚由何而来。

\section{传统方法为何失效}

表~20.1 显示的是来自一项存在处理—混杂反馈的序贯随机化试验的数据，如图~20.3 的因果图所示。尽管没有未测量变量 $U_1$（免疫抑制水平）的数据，三个可识别性条件仍然全部成立：只要掌握混杂因子 $L_1$ 的数据，就不需要 $U_1$。因此，如第 19 章所述，我们应当能够正确估计涉及任何静态或动态处理策略的因果效应。

然而，上一节的分析无论是否校正 $L_1$，都没有给出正确答案。

问题在于我们没用对校正混杂的方法。分层是校正混杂的常用方法，却无法处理处理—混杂反馈。分层意味着在由混杂因子——本例中为 $L_1$——定义的研究总体子集（层）中估计处理与结局之间的关联。由于变量 $L_1$ 只能取两个值——CD4 细胞计数低时为 1，否则为 0——本例中有两个这样的层。为估计 $L_1=l$ 者的因果效应，我们选取（即以 $L_1=l$ 为条件、或按 $L_1=l$ 分层）总体中取值为 $L_1=l$ 的子集。

\begin{figure}[htbp]
\centering
\begin{tikzpicture}[causal,every node/.style={transform shape},scale=1.15]
  \node[nd] (A0) at (0,0) {$A_0$};
  \node[bx] (L1) at (2.4,0) {$L_1$};
  \node[nd] (A1) at (4.4,0) {$A_1$};
  \node[nd] (Y)  at (6.1,0) {$Y$};
  \node[nd] (U1) at (2.4,-1.5) {$U_1$};
  \draw[ed] (A0) -- (L1);
  \draw[ed] (L1) -- (A1);
  \draw[ed] (U1) -- (L1);
  \draw[ed] (U1) .. controls (4.8,-1.15) and (5.8,-0.75) .. (Y);
\end{tikzpicture}
\caption{以 $L_1$（方框）为条件，打开了路径 $A_0\to L_1\leftarrow U_1\to Y$}
\label{fig:20.5}
\end{figure}

然而，当关联度量是在由先前处理 $A_0$ 所致的变量 $L_1$ 的各水平内计算时，分层可能产生意想不到的后果。确实，图~20.5 表明，以 $L_1$——一个对撞节点——为条件，打开了路径 $A_0\to L_1\leftarrow U_1\to Y$。也就是说，分层在 $L_1$ 的各水平内，诱导出第 0 期的处理 $A_0$ 与未测量变量 $U_1$ 之间的非因果关联，从而也诱导出 $A_0$ 与结局 $Y$ 之间的非因果关联。在 CD4 计数低（$L_1=1$）者中，受处理（$A_0=1$）成了严重免疫抑制（$U_1$ 取值高）的标志；在 CD4 水平高（$L_1=0$）者中，未受处理（$A_0=0$）成了免疫抑制较轻（$U_1$ 取值低）的标志。于是，分层的副作用就是在处理 $A_0$ 与结局 $Y$ 之间诱导出关联。

换言之，分层消除了 $A_1$ 的混杂，代价是为 $A_0$ 引入了选择偏倚。即使处理在任何时点对任何个体的结局都没有效应（如本例），关联差
\[
\E[Y\given A_0=1,L_1=l,A_1=1]-\E[Y\given A_0=0,L_1=l,A_1=0]
\]
也可能不为 0。这一偏倚源于按一个受（时变处理的成分 $A_0$ 所）影响的变量 $L_1$ 来选取研究总体的子集。净偏倚取决于被消除的混杂与被引入的选择偏倚的相对大小。

\begin{figure}[htbp]
\centering
\begin{tikzpicture}[causal,every node/.style={transform shape},scale=1.15]
  \node[nd] (A0) at (0,0) {$A_0$};
  \node[bx] (L1) at (2.5,0) {$L_1$};
  \node[nd] (A1) at (4.5,0) {$A_1$};
  \node[nd] (Y)  at (6.2,0) {$Y$};
  \node[nd] (W0) at (0,-1.5) {$W_0$};
  \node[nd] (U1) at (2.5,-1.5) {$U_1$};
  \draw[ed] (W0) -- (A0);
  \draw[ed] (W0) -- (L1);
  \draw[ed] (L1) -- (A1);
  \draw[ed] (U1) -- (L1);
  \draw[ed] (U1) .. controls (4.9,-1.15) and (5.9,-0.75) .. (Y);
\end{tikzpicture}
\caption{以对撞节点 $L_1$（方框）为条件，打开了路径 $A_0\leftarrow W_0\to L_1\leftarrow U_1\to Y$}
\label{fig:20.6}
\end{figure}

严格说来，传统方法的偏倚不仅出现在混杂因子受先前处理影响时（在随机化实验或观察性研究中），也出现在混杂因子与先前处理共享未测量原因 $W_0$ 时（在观察性研究中）。在图~20.6 所刻画的观察性研究中，以对撞节点 $L_1$ 为条件打开了路径 $A_0\leftarrow W_0\to L_1\leftarrow U_1\to Y$。因此，我们把图~20.3 与图~20.4 这两种（无法用观测数据区分）情形都称为处理—混杂反馈的例子。

\begin{figure}[htbp]
\centering
\begin{tikzpicture}[causal,every node/.style={transform shape},scale=1.15]
  \node[nd] (A0) at (0,0) {$A_0$};
  \node[bx] (L1) at (2.4,0) {$L_1$};
  \node[nd] (A1) at (4.4,0) {$A_1$};
  \node[nd] (Y)  at (6.1,0) {$Y$};
  \node[nd] (U1) at (2.0,-1.5) {$U_1$};
  \draw[ed] (A0) -- (L1);
  \draw[ed] (L1) -- (A1);
  \draw[ed] (L1) .. controls (4.3,1.05) and (6.2,0.95) .. (Y);
  \draw[ed] (U1) -- (L1);
  \draw[ed] (U1) .. controls (4.9,-1.3) and (5.9,-0.8) .. (Y);
\end{tikzpicture}
\caption{混杂因子 $L_1$ 位于因果路径 $A_0\to L_1\to Y$ 上}
\label{fig:20.7}
\end{figure}

我们用来描述传统方法偏倚的因果图都很简单，只刻画了处理对结局没有因果效应的情形。然而，存在处理—混杂反馈时，以混杂因子为条件在处理效应非零时同样会诱导偏倚，如图~20.7。从 $A_0$、$A_1$ 或 $L_1$ 指向 $Y$ 的箭头是否存在，并不改变这一事实：以 $L_1$ 为条件会在 $A_0$ 与 $Y$ 之间造成没有因果解释的关联（另见精要~20.2）。此外，我们的因果图只有两个时点、节点数目也有限，但传统方法的偏倚同样会出现在包含多个时点与变量的高维数据中。事实上，多个时点上存在受先前处理影响的时变混杂因子，会增大出现大幅偏倚的可能。

\begin{finepoint}[因果路径上的混杂因子]
以受先前处理影响的混杂因子 $L_1$ 为条件，即使该混杂因子不在处理与结局之间的因果路径上，也可能造成选择偏倚。事实上，图~20.5 与图~20.6 中并不存在这样的因果路径。

另一方面，在图~20.7 中，后续处理 $A_1$ 的混杂因子 $L_1$ 位于从先前处理 $A_0$ 到结局 $Y$ 的因果路径上，即路径 $A_0\to L_1\to Y$。若图~20.7 中 $U_1$ 不是 $L_1$ 与 $Y$ 的共同原因（即若不存在选择偏倚），则 $L_1$ 各层内 $A$ 与 $Y$ 的关联将是 $A_0$ 不经 $L_1$ 的直接效应的无偏估计，但仍不会是 $\bar A$ 之总效应的无偏估计，因为经 $L_1$ 中介的 $A_0$ 效应没有被计入。

人们有时说，处理与结局之间因果路径上的变量不能视为混杂因子，因为校正这些变量会使效应估计有偏。然而，对时变处理而言，这种对混杂因子的刻画并不准确。图~20.7 表明，后续处理 $A_1$ 的混杂因子可以位于过去处理 $A_0$ 与结局之间的因果路径上。至于校正因果路径上的混杂因子是否会给处理策略的效应带来偏倚，则取决于所选的校正方法。分层确实会引入偏倚；g 方法则不会。
\end{finepoint}

一般而言，只有当处理各成分 $A_k$ 的联合效应能够同时且无偏地估计时，才可能有效估计处理策略的效应。正如刚才所见，即使掌握了全部时变混杂因子的数据，用分层也可能无法做到。

\section{传统方法为何无法修补}

我们已经说明，存在处理—混杂反馈时，分层不能用作混杂校正方法。那么其他传统方法呢？例如，我们本可以用参数结局回归来校正混杂，而不必用非参数分层。结局回归能在简单分层失败的地方取得成功吗？

这个问题对高维数据的情形尤为重要，因为在高维情形下，我们无法像上一节那样做简单的分层分析。考虑按图~20.5 生成的数据。处理 $A_k$ 发生在 $k=0,1$ 两个月，这意味着只有 $2^2=4$ 种静态处理策略 $\bar a$。但当处理 $A_k$ 发生在多个时点 $k=0,1\ldots K$ 时，我们无法把所有处理取值组合列成一张表。若 $K$ 的量级为 100（实践中并不少见），就有 $2^{100}$ 种静态处理策略 $\bar a$，这个惊人的数目远超任何研究的样本量。\mnote{若每个时点 $k$ 都测量了多个混杂因子 $L_k$，数据组合的数目还要大得多。}若把动态策略也考虑在内，处理策略的总数还要大得多。

正如我们从第 11 章起一直主张的，当可能的处理策略 $\bar a$ 很多时，我们需要建模来估计涉及 $\E[Y^{\bar a}]$ 的平均因果效应。为此，需要为处理史 $\bar a$ 对平均结局 $Y$ 的效应假定一个剂量—反应函数。一种可能是假定处理策略 $\bar a$ 的效应随各策略下的累积处理量线性增加。在这一假定下，所有恰好分配三个月治疗的策略效应相同，无论这三个月治疗出现在哪一段——随访的头三个月、随访的最后三个月，等等。建模的代价是：剂量—反应函数可能误设，从而给估计的有效性带来又一个威胁。

然而，付出这一代价，并不能换来对传统方法失效的防护。存在处理—混杂反馈时，回归建模不可能消除常规分层方法的偏倚，因为回归本身就是一种常规的基于分层的方法。例如，设我们有按图~20.5 生成的数据。定义累积处理 $\mathrm{cum}(\bar A)=A_0+A_1$，它可取 3 个值：0（个体两个时点都未受处理）、1（受试者只在第 1 期或只在第 2 期受处理）、2（受试者两个时点都受处理）。所关心的处理策略于是可表示为“始终治疗”$\mathrm{cum}(\bar a)=2$ 与“从不治疗”$\mathrm{cum}(\bar a)=0$，平均因果效应为
\[
\E\bigl[Y^{\mathrm{cum}(\bar a)=2}\bigr]-\E\bigl[Y^{\mathrm{cum}(\bar a)=0}\bigr]
\]
同样，任何有效方法都应当估计出这一差值为 0。

在平均结局 $\E[Y\given \bar A,L_1]$ 线性依赖于协变量 $\mathrm{cum}(\bar A)$ 的假定下，我们可以拟合结局回归模型
\[
\E[Y\given \bar A,L_1]=\theta_0+\theta_1\,\mathrm{cum}(\bar A)+\theta_2L_1
\]

关联差
\[
\E\bigl[Y\given \mathrm{cum}(\bar A)=2,L_1\bigr]-\E\bigl[Y\given \mathrm{cum}(\bar A)=0,L_1\bigr]
\]
等于 $\theta_1\times 2$。（模型正确地假定该差值在层 $L_1=1$ 与 $L_1=0$ 中相同。）因此有人可能想把 $\theta_1\times 2$ 解释为“始终治疗”相对于“从不治疗”在协变量 $L_1$ 各水平内的平均因果效应。但这样的因果解释没有依据，因为如图~20.5 所示，以 $L_1$ 为条件会在 $A_0$——处理的成分 $\mathrm{cum}(\bar A)$ 之一——与结局 $Y$ 之间诱导出关联。这意味着，即使真实因果效应为零、且 $\E[Y\given \bar A,L_1]$ 的回归模型正确，$\theta_1$（因而均值之关联差）也不为零。

类似的论证也适用于匹配。存在处理—混杂反馈时，要适当校正时变混杂因子，需要 g 方法。

\section{校正过去处理}

在讨论 g 方法之前，还有一点。为简单起见，我们迄今都在简化的因果图下描述处理—混杂反馈，其中过去处理不直接影响后续处理。也就是说，图~20.3 与图~20.4 的因果图没有包含从 $A_0$ 指向 $A_1$ 的箭头。现在我们考虑更一般的情形：过去处理可能直接影响后续处理。

例如，设我们 HIV 研究中的医生在决定第 $k$ 期是否开出处理 $A_k$ 时，使用了过去处理史 $\bar A_{k-1}$ 的信息。为表示这一情形，我们在图~20.3 与图~20.4 的因果图中加入从 $A_0$ 指向 $A_1$ 的箭头，如图~20.8 与图~20.9 所示。

\begin{figure}[htbp]
\centering
\begin{tikzpicture}[causal,every node/.style={transform shape},scale=1.15]
  \node[nd] (A0) at (0,0) {$A_0$};
  \node[nd] (L1) at (2.6,0) {$L_1$};
  \node[nd] (A1) at (4.5,0) {$A_1$};
  \node[nd] (Y)  at (6.2,0) {$Y$};
  \node[nd] (U1) at (2.6,-1.5) {$U_1$};
  \draw[ed] (A0) -- (L1);
  \draw[ed] (A0) to[out=48,in=132] (A1);
  \draw[ed] (L1) -- (A1);
  \draw[ed] (U1) -- (L1);
  \draw[ed] (U1) .. controls (4.9,-1.15) and (5.9,-0.75) .. (Y);
\end{tikzpicture}
\caption{图~20.3 加入从 $A_0$ 指向 $A_1$ 的箭头}
\label{fig:20.8}
\end{figure}

\begin{figure}[htbp]
\centering
\begin{tikzpicture}[causal,every node/.style={transform shape},scale=1.15]
  \node[nd] (A0) at (0,0) {$A_0$};
  \node[nd] (L1) at (2.7,0) {$L_1$};
  \node[nd] (A1) at (4.6,0) {$A_1$};
  \node[nd] (Y)  at (6.3,0) {$Y$};
  \node[nd] (W0) at (0,-1.5) {$W_0$};
  \node[nd] (U1) at (2.7,-1.5) {$U_1$};
  \draw[ed] (W0) -- (A0);
  \draw[ed] (W0) -- (L1);
  \draw[ed] (A0) to[out=48,in=132] (A1);
  \draw[ed] (L1) -- (A1);
  \draw[ed] (U1) -- (L1);
  \draw[ed] (U1) .. controls (5.0,-1.15) and (6.0,-0.75) .. (Y);
\end{tikzpicture}
\caption{图~20.4 加入从 $A_0$ 指向 $A_1$ 的箭头}
\label{fig:20.9}
\end{figure}

图~20.8 与图~20.9 的因果图表明，存在处理—混杂反馈时，以 $L_1$ 为条件不足以阻断处理 $A_1$ 与结局 $Y$ 之间的全部后门路径。确实，以 $L_1$ 为条件打开了图~20.8 中的路径 $A_1\leftarrow A_0\to L_1\leftarrow U_1\to Y$，以及图~20.9 中的路径 $A_1\leftarrow A_0\leftarrow W_0\to L_1\leftarrow U_1\to Y$。当然，无论是否存在处理—混杂反馈，只要过去处理对结局有非零效应（如图~20.10 的因果图），就始终需要以过去处理史为条件。在该图下，处理 $A_0$ 是处理 $A_1$ 之效应的混杂因子。

\begin{figure}[htbp]
\centering
\begin{tikzpicture}[causal,every node/.style={transform shape},scale=1.15]
  \node[nd] (A0) at (0,0) {$A_0$};
  \node[nd] (L1) at (2.6,0) {$L_1$};
  \node[nd] (A1) at (4.5,0) {$A_1$};
  \node[nd] (Y)  at (6.2,0) {$Y$};
  \node[nd] (U1) at (2.3,-1.5) {$U_1$};
  \draw[ed] (A0) -- (L1);
  \draw[ed] (A0) .. controls (3.0,0.95) and (4.5,0.95) .. (A1);
  \draw[ed] (A0) .. controls (3.2,1.90) and (5.6,1.80) .. (Y);
  \draw[ed] (L1) -- (A1);
  \draw[ed] (U1) -- (L1);
  \draw[ed] (U1) .. controls (5.0,-1.3) and (6.0,-0.8) .. (Y);
\end{tikzpicture}
\caption{过去处理 $A_0$ 直接影响 $Y$ 时，$A_0$ 成为 $A_1$ 之效应的混杂因子}
\label{fig:20.10}
\end{figure}

因此，第 $k$ 期的序贯可交换性一般要求以 $k$ 之前的处理史 $\bar A_{k-1}$ 为条件；只以协变量 $L$ 为条件是不够的。这就是为什么在本章与上一章中，所有表示序贯可交换性的条件独立陈述都以处理史为条件。

过去处理在估计时点固定处理的效应时同样重要。设我们关心的是估计时点固定处理 $A_1$ 对 $Y$ 的效应——而不是涉及 $A_0$ 与 $A_1$ 两者的处理策略的效应。（$A_1$ 的效应有时被称为时变处理 $\bar A$ 的短期效应。）那么，不校正过去处理 $A_0$，在存在处理—混杂反馈时一般会导致选择偏倚，在过去处理 $A_0$ 直接影响结局 $Y$ 时则会导致混杂。换言之，即使处理 $A_1$ 对任何个体的结局 $Y$ 都没有效应（如图~20.8—20.10），差值 $\E[Y\given A_1=1,L_1]-\E[Y\given A_1=0,L_1]$ 也不会为零。在实践中，对时点固定处理做因果推断时，比较处理现用者（$A_1=1$）与非使用者（$A_1=0$）的分析可能产生偏倚。为避免偏倚，可以校正先前处理史，或把分析限制在具有特定处理史的个体中。这正是时点固定处理的“新使用者设计”背后的思想：把分析限制在过去未使用过处理的个体中。\mnote{若能正确校正过去处理，分析就不必限于新使用者。}

当过去处理被误测时，校正过去处理的要求还有额外的偏倚含义。如 9.3 节所述，误测的混杂因子可能使效应估计向上或向下偏倚。在我们的 HIV 例子中，设研究者无法获得研究参与者的病历。为确定先前处理，研究者不得不通过问卷询问参与者。由于并非所有参与者都能准确回忆自己的处理史，处理 $A_0$ 的测量带有误差。研究者手头有的是误测变量 $A_0^{*}$，而不是变量 $A_0$。为在图~20.8—20.10 中刻画这一情形，我们加入从真实处理 $A_0$ 指向误测处理 $A_0^{*}$ 的箭头；这表明以 $A_0^{*}$ 为条件无法阻断 $A_1$ 与 $Y$ 之间经过 $A_0$ 的偏倚路径。于是，尽管 $A_1$ 对 $Y$ 没有效应，研究者在校正 $A_0^{*}$ 与 $L_1$ 之后仍会得出 $A_1$ 与 $Y$ 之间存在关联的结论。

因此，当处理是时变的时候，我们发现，与一种流传甚广的看法相反，处理的误测——即使测量误差是独立且非差异的——也可能在零假设下造成偏倚。这一偏倚之所以出现，是因为即使过去处理对结局没有因果效应，它仍是后续处理之效应的混杂因子。此外，在备择假设下，这种不完美的偏倚校正可能导致效应估计被夸大。\mnote{Robins (1987) 表明，随机误测的处理可能导致远离零值的偏倚。}
